\chapter*{CAPÍTULO 3. APLICACIÓN DE SOCIALSCRUM EN EQUIPOS DE DESARROLLO DE SOFTWARE ÁGIL}
\addcontentsline{toc}{chapter}{CAPÍTULO 3. APLICACIÓN DE SOCIALSCRUM EN EQUIPOS DE DESARROLLO DE SOFTWARE ÁGIL}
\setcounter{chapter}{3}
\setcounter{section}{0}
\setcounter{subsection}{0}
\setcounter{subsubsection}{0}
\renewcommand{\thesection}{\thechapter.\arabic{section}}
\renewcommand{\thesubsection}{\thesection.\arabic{subsection}}
\renewcommand{\thesubsubsection}{\thesubsection.\arabic{subsubsection}}

Este capítulo describe la aplicación de SocialScrum, un proceso ágil adaptado para gestionar la deuda social en equipos de desarrollo de software. Se centra en la implementación práctica de los eventos y artefactos modificados, así como en el rol del Social Guide, para abordar problemas interpersonales y mejorar el bienestar del equipo.